\documentclass{article}
\usepackage{amsmath}
\usepackage{cancel}


\title{Taylor Series in One Dimension}
\begin{document}
We would like to be able to approximate the value of a function in a small
neighborhood of a point by a series of polynomials. This is useful for many reasons.
Polynomials are far easier to work with than almost every other function. They
require only addition and multiplication to define and are very simple to
differentiate and integrate. Expanding a function into a polynomial series makes
it easy to take nasty limits and can even aid in solving stubborn differential
equations. 


To begin let's consider a function $f$. We want to find the coefficients so that
\begin{equation}
  f(x) = \sum\limits_{n=0}^\infty c_n(x-a)^n
\end{equation}
Expanding this series for the first few terms, we have
\begin{equation}
  f(x) = c_0 + c_1(x-a) + c_2(x-a)^2 + c_3(x-a)^3 + ...
\end{equation}
Notice how we have $(x-a)^n$ all over the place. If we let $x=a$, then we find
\begin{align}
  f(a) &= c_0 + 0 + 0 + 0 + ...  = c_0   \\
  \Rightarrow c_0 &= f(a)
\end{align}
Now if we take one derivative, we have
\begin{align}
  f'(x) &= c_1 + 2c_2(x-a) + 3c_3(x-a)^2 + ... \\
  f'(a) &= c_1 + 0 + 0 + ... \\
  \Rightarrow c_1 &= f'(a)
\end{align}
Taking another derivative, we have
\begin{align}
  f''(x)&= 2c_2  +6c_3(x-a) + ... \\
  f''(a)&= 2c_2 + 0 + ... \\
  \Rightarrow c_2 &= \frac{1}{2}f''(a)
\end{align}
and again...
\begin{align}
  f^{(3)}(x) &= 6c_3 + ... \\
  f^{(3)}(a) &= 6c_3 + 0 + ... \\
  \Rightarrow c_3 &= \frac{1}{6}f^{(3)}(a)
\end{align}

At this point, we can hopefully see the pattern. We have
\begin{align}
  c_0 &= f(a) = \frac{1}{0!}f^{(0)}(a) \\
  c_1 &= f'(a) = \frac{1}{1!}f^{(1)}(a) \\
  c_2 &= \frac{1}{2}f''(a) = \frac{1}{2!}f^{(2)}(a) \\
  c_3 &= \frac{1}{6}f^{(3)}(a) = \frac{1}{3!}f^{(3)}(a) \\
  \vdots &= \vdots \\
  c_n &= \frac{1}{n!}f^{(n)}(a)
\end{align}
Putting this all together, we have found that the Taylor series for a suitably
differentiable function $f$ about the point $x=a$ is given by
\begin{equation}
  f(x) = \sum\limits_{n=0}^\infty\frac{f^{(n)}(a)}{n!}(x-a)^n
\end{equation}

Let's try and visualize these approximations using Sage for
\begin{equation}
  f(x) = \sin(x)
\end{equation}
with $a = \pi/6$. 

%% taylor-series-1d.sage %%


\end{document}
